\documentclass[11pt]{article}

\usepackage[margin=0.95in]{geometry}
\usepackage[T1]{fontenc}
\usepackage[utf8]{inputenc}
\usepackage{lmodern}
\usepackage{microtype}
\usepackage[hidelinks]{hyperref}
\usepackage{booktabs}
\usepackage{xcolor}
\usepackage{amsmath}
\usepackage{array}

\definecolor{accent}{RGB}{18,54,92}
\setlength{\parindent}{0pt}
\setlength{\parskip}{8pt}

\title{\textbf{A Full Explanation of the Spectral Signal Extraction for Statistical Arbitrage Project}}
\author{}
\date{\today}

\begin{document}
\maketitle
\thispagestyle{empty}

\section*{Introduction}

This project studies a specific problem in cross-sectional statistical
arbitrage. The question is not whether one can write down a new theorem in
random matrix theory, and it is not whether one can claim a magical source of
alpha. The question is narrower and more practical: if a stat-arb signal is
built from residual mean reversion, does the \emph{quality of residual
construction} improve when we use a noise-aware spectral filtering rule rather
than a naive fixed-rank PCA baseline?

That is the heart of the project. The strategy is only the testing ground. The
real object of study is the residualization step.

\section*{What Problem the Project Is Solving}

Cross-sectional equity stat-arb usually begins with the idea that stocks move
for two broad reasons. Some part of a stock's return is common and systematic:
the market went up, a sector sold off, a style factor moved, or a group of
names shared the same broad exposure. Another part of the return is
idiosyncratic: the stock moved more or less than what its common exposures would
predict. That leftover component is the residual.

Residual mean reversion strategies try to trade that leftover component. If a
stock looks unusually rich relative to its common factor structure, perhaps it
should be shorted. If it looks unusually cheap, perhaps it should be bought.
The entire success of such a strategy depends on how well one separates the
common component from the residual component. If that separation is poor, then
the signal is contaminated from the start.

This is where covariance estimation becomes important. To estimate the common
structure, one often uses PCA on a rolling return panel. But sample covariance
matrices are noisy in finite samples. If the number of assets is not very small
relative to the rolling window, empirical eigenvalues and eigenvectors can be
unstable. A naive PCA decomposition may treat noise-dominated modes as real
structure, or remove the wrong directions entirely. The project asks whether
Marchenko--Pastur-style filtering can reduce that problem enough to matter in a
trading workflow.

\section*{Background Linear Algebra Intuition}

One useful way to motivate the project is through the classical
eigenportfolio viewpoint. If \(w\) is a portfolio weight vector and
\(\Sigma\) is the covariance matrix of asset returns, then the quadratic form
\[
w^\top \Sigma w
\]
measures the portfolio's variance. In that language, eigenvectors of
\(\Sigma\) can be interpreted as orthogonal risk directions, and eigenvalues
measure how much variance those directions explain.

That perspective matters here for two reasons. First, it gives a concrete
meaning to the spectrum of the covariance matrix: the largest modes are not
abstract linear algebra objects but the dominant sources of common movement in
the panel. Second, in broad equity data the top eigenvalue is often
market-like, while the next few modes often capture sector or style structure.
So even before we talk about RMT, linear algebra already tells us why a
residualization problem exists at all: if those common modes are not removed
carefully, the supposed stock-specific signal is polluted from the start.

The key difference between this project and the older eigenportfolio story is
the target. That classical setup asks which eigenvector one might hold as a
portfolio. This project instead asks which spectral directions should be
treated as common structure and stripped out before the remainder is traded as
a residual mean-reversion signal. So the eigendecomposition is serving signal
extraction rather than buy-and-hold allocation.

\section*{The Core Research Question}

The cleanest version of the question is:

\begin{quote}
If we hold the universe, dates, signal construction, neutrality constraints,
turnover rules, and transaction-cost assumptions fixed, and change only the
residualization method, do we get a stronger residual mean-reversion sleeve from
RMT-filtered PCA than from raw fixed-rank PCA?
\end{quote}

That is a meaningful applied research question because it isolates the effect of
the spectral estimation step itself.

\section*{What Data the Public Repository Uses}

The public repository uses a structured benchmark panel rather than a
redistributed historical equity universe. This is documented explicitly in
\texttt{src/sample\_data.py}. The synthetic panel contains 24 equities across 6
sectors and spans business dates from \texttt{2019-01-02} through
\texttt{2025-06-30}. The returns are generated from a combination of:

\begin{itemize}
    \item a market factor
    \item a style factor
    \item one sector factor per sector
    \item a sector-level mean-reversion state
    \item asset-specific residual noise
    \item a common cross-sectional residual shock
\end{itemize}

This design is not trying to pretend that the public data is live market data.
Its purpose is different. It creates a transparent finite-sample covariance
estimation problem in which naive PCA can overfit unstable spectral directions,
while still leaving a residual structure that can be tested systematically. In
other words, the synthetic benchmark is a methodology testbed, not a production
claim.

\section*{How the Original Historical Project Relates to the Public Repo}

The original Georgia Tech project used a broader historical equity panel that is
not redistributed. The preserved summary documents the following approximate
historical comparison:

\begin{center}
\small
\begin{tabular}{@{}>{\raggedright\arraybackslash}p{0.34\textwidth} >{\raggedleft\arraybackslash}p{0.14\textwidth} >{\raggedright\arraybackslash}p{0.34\textwidth}@{}}
\toprule
\textbf{Method} & \textbf{Approximate net Sharpe} & \textbf{Notes} \\
\midrule
Raw PCA residualization & 0.8 & Historical walk-forward comparison baseline \\
RMT-filtered residualization & 1.3 & Historical walk-forward comparison winner \\
\bottomrule
\end{tabular}
\normalsize
\end{center}

The preserved materials also report max drawdown around \(8\%\) and benchmark
beta with \(|\beta_{\mathrm{SPY}}| < 0.05\). The public repository therefore has
two roles:

\begin{itemize}
    \item it is a reproducible benchmark implementation of the methodology
    \item it is a preserved documentation layer for the original historical project
\end{itemize}

It should not be read as claiming that the synthetic panel reproduces the exact
historical run.

\section*{How the Spectral Methods Differ}

The benchmark compares several methods, but the central comparison is between
raw PCA residualization and Marchenko--Pastur-filtered residualization.

Raw PCA residualization works as follows. Take a rolling return window, compute
the correlation structure, perform an eigendecomposition, select a fixed number
of leading components, reconstruct the common component using those directions,
and subtract it from the standardized return panel. The result is a matrix of
residual returns that are supposed to represent stock-specific deviations after
common structure is removed.

The weakness is that fixed-rank PCA assumes in advance how many factors should
be treated as real. It also assumes that the top empirical modes are stable
enough to be trusted. In finite samples, both assumptions can fail.

The RMT version uses the Marchenko--Pastur upper edge as a filtering rule. The
empirical spectrum of the sample correlation matrix is compared against the
noise threshold implied by the dimension-to-sample-size ratio. Only eigenvalues
above that edge are treated as significant common structure. Those selected
directions are then used to reconstruct the common component, and the remainder
is treated as the residual panel. In the repo, that logic lives in
\texttt{src/spectral.py}.

In the current implementation, the rolling spectral step follows a classical
correlation-filtering convention: returns inside the window are transformed
with \(\log(1+r)\), standardized asset by asset, and then mapped into a sample
correlation matrix before eigendecomposition. The backtest still evaluates
arithmetic close-to-close returns after costs; only the spectral estimation
step uses the log-standardized input.

So the new information provided by the RMT step is not a prediction about future
returns by itself. The new information is an estimate of which spectral modes
are more likely to reflect real common structure rather than sampling noise.
That changes the quality of the residuals that feed the downstream trading
signal.

\section*{How the Trading Signal Is Built}

Once residuals are computed, the project turns them into a cross-sectional
mean-reversion signal. The residual window is standardized into a recent z-score
for each asset. The sign convention is mean-reverting: unusually positive
residuals are treated as candidates to short, and unusually negative residuals
as candidates to buy. Only the strongest names are selected on each side, so the
strategy is trading relative dislocations rather than trying to hold every name.

The project then builds a long/short portfolio with several layers of control.
Weights are made approximately dollar-neutral so that the book is not simply a
market bet. They are beta-neutralized with respect to the benchmark, which in
the public repo is analogous to SPY beta control. Gross exposure is normalized,
and per-position caps are enforced so that the strategy does not concentrate
heavily in a few names. Some sleeves also blend toward prior weights, which is a
practical attempt to control turnover.

\section*{How the Backtest Works}

The backtest is intentionally conservative. Portfolio weights chosen at time
\(t\) are applied to the next period's returns, not to same-day returns. That
shift matters because it avoids accidental look-ahead bias. Turnover is computed
from changes in target weights, and one-way transaction cost is charged on that
turnover. The public benchmark uses a 5 bps one-way cost assumption as its base
case.

So by the time the strategy is evaluated, the comparison between raw PCA and RMT
is no longer just a comparison between two residual matrices. It is a comparison
between two fully implemented, cost-aware, beta-controlled, dollar-neutral
stat-arb sleeves.

\section*{How the Public Benchmark Is Split}

The public benchmark uses a simple train / validation / holdout structure:

\begin{center}
\begin{tabular}{lll}
\toprule
\textbf{Split} & \textbf{Dates} & \textbf{Purpose} \\
\midrule
Train & 2019-01-02 to 2021-12-31 & First-pass design \\
Validation & 2022-01-03 to 2022-12-30 & Selection of the final spectral sleeve \\
Holdout & 2023-01-02 to 2025-06-30 & Final public out-of-sample evaluation \\
\bottomrule
\end{tabular}
\end{center}

The public selection is validation-only. Among the spectral sleeves, the chosen
final public sleeve is \texttt{rmt\_filtered\_residual}. The public repo also
publishes that same strategy again under the name
\texttt{final\_research\_portfolio} so that review-facing summary tables are
easy to read.

\section*{What the Public Results Show}

The headline public benchmark table is:

\begin{center}
\small
\begin{tabular}{@{}>{\raggedright\arraybackslash}p{0.30\textwidth} >{\raggedleft\arraybackslash}p{0.11\textwidth} >{\raggedleft\arraybackslash}p{0.12\textwidth} >{\raggedleft\arraybackslash}p{0.12\textwidth} >{\raggedleft\arraybackslash}p{0.10\textwidth} >{\raggedleft\arraybackslash}p{0.11\textwidth}@{}}
\toprule
\textbf{Strategy} & \shortstack{\textbf{Val}\\\textbf{Sharpe}} & \shortstack{\textbf{Holdout}\\\textbf{Sharpe}} & \shortstack{\textbf{Holdout}\\\textbf{Max DD}} & \shortstack{\textbf{Holdout}\\\textbf{Beta}} & \shortstack{\textbf{Holdout}\\\textbf{Turnover}} \\
\midrule
Raw PCA Residual & -0.87 & -0.21 & -5.1\% & 0.011 & 1.52 \\
RMT-Filtered Residual & 1.04 & 1.73 & -2.7\% & -0.004 & 1.54 \\
Final Research Portfolio & 1.04 & 1.73 & -2.7\% & -0.004 & 1.54 \\
\bottomrule
\end{tabular}
\normalsize
\end{center}

The direction of the result is clear. In the public benchmark, the raw PCA
sleeve performs poorly out of sample, while the RMT-filtered sleeve remains
strongly positive after costs and has better drawdown and beta behavior.

The most important interpretive point is that the intended difference between
those sleeves is the residualization step. Everything else in the comparison is
held as constant as possible. This is why the result is meaningful as a research
finding: it suggests that covariance denoising is not cosmetic in a residual
mean-reversion strategy. It can materially change the quality of the downstream
signal.

\section*{What the Cost Sensitivity Tells Us}

The public cost sensitivity table is severe:

\begin{center}
\begin{tabular}{lrrr}
\toprule
\textbf{Cost} & \textbf{Holdout Sharpe} & \textbf{Holdout Ann Return} & \textbf{Holdout Max DD} \\
\midrule
1 bps & 5.63 & 24.9\% & -1.3\% \\
5 bps & 1.73 & 6.9\% & -2.7\% \\
10 bps & -3.36 & -12.0\% & -27.9\% \\
20 bps & -14.26 & -40.4\% & -73.6\% \\
\bottomrule
\end{tabular}
\end{center}

This is not something to hide. It is one of the most informative parts of the
project. Residual mean-reversion strategies often trade frequently, and frequent
trading makes them fragile to execution costs. So the public benchmark should
not be read as proof of a live deployable edge. It should be read as proof that
the residual-denoising question matters under a stylized but transparent cost
model. The strategy's next research problem is clearly turnover reduction and
better execution realism.

\section*{What the Robustness Checks Say}

The parameter-sensitivity results are also informative. When the rolling window
is changed by \(\pm 20\%\), the RMT-filtered sleeve remains positive and fairly
stable on holdout, while the raw PCA sleeve remains much weaker. That matters
because it suggests the RMT advantage is not coming purely from a knife-edge
parameter choice.

The regime summaries show that the RMT sleeve remains particularly strong in the
public benchmark's bull/high-volatility and event-window recovery periods, while
also avoiding the drawdown severity of the raw PCA version. So the improvement
is not merely one lucky average over the whole sample; it also appears in the
way the sleeve behaves across different structural environments.

\section*{What We Can Infer from the Project}

The project supports a few strong conclusions. First, the quality of
residualization is part of the signal model itself, not just a preprocessing
detail. If common-structure estimation is noisy, then the residual mean-reversion
signal can be badly degraded before the trading logic even begins. Second,
Marchenko--Pastur-style filtering can improve that residualization step in a
controlled high-dimensional setting. Third, the improvement can survive the
addition of dollar neutrality, beta control, and transaction costs. Those are
meaningful research findings even if they are not equivalent to a production
trading claim.

The project also suggests a more conceptual point: random matrix theory is
useful here not because it ``predicts returns'' directly, but because it gives a
better estimate of which covariance directions should be treated as structure and
which should be treated as noise. That is the real contribution of the theory in
this context.

\section*{What the Project Does Not Prove}

The public benchmark does not prove a live tradable stat-arb edge. It does not
fully model borrow fees, short frictions, market impact, or realistic universe
evolution. It does not replace a point-in-time historical U.S. equity panel. It
also does not, by itself, fully prove the original Georgia Tech historical
performance claim. The historical claim is preserved through summary artifacts,
not through a committed raw output table from the original private run.

That limitation should be stated clearly, because the project is stronger when
it is interpreted honestly. It is a serious methodological experiment on
residual-denoising inside a stat-arb workflow. It is not a production-ready
trading system.

\section*{Conclusion}

The best way to understand the statistical arbitrage project is as a carefully
structured test of whether better spectral denoising produces better residual
signals. The public repository shows the full workflow: generate or load a
return panel, estimate covariance structure, residualize via raw PCA or RMT,
turn residuals into a long/short mean-reversion signal, impose neutrality and
beta controls, charge transaction costs, and compare the resulting sleeves on a
validation / holdout split. The original Georgia Tech project suggests that this
idea also improved a broader historical run from about 0.8 Sharpe to about 1.3
Sharpe under comparable constraints. Even where the public benchmark remains
stylized, the project succeeds in showing that the residualization step itself
is an intellectually serious and practically important part of a cross-sectional
stat-arb pipeline.

\end{document}
